\documentclass{article}%
\usepackage[T1]{fontenc}%
\usepackage[utf8]{inputenc}%
\usepackage{lmodern}%
\usepackage{textcomp}%
\usepackage{lastpage}%
\usepackage[backend=biber]{biblatex}%
\usepackage{fancyhdr}%
\usepackage{lastpage}%
\usepackage{ragged2e}%
\usepackage{pdfpages}%
\usepackage{hyperref}%
\usepackage{booktabs}%
\usepackage{float}%
\usepackage{threeparttable}%
\usepackage{amssymb}%
\usepackage{amsmath}%
\usepackage{adjustbox}%
\usepackage{geometry}%
\usepackage{fancyhdr}%
%
\geometry{left=1in, right=1in, top=1in, bottom=1in}%
\addbibresource{references.bib}%
\hypersetup{colorlinks=true, linkcolor=blue, urlcolor=blue}%
\fancypagestyle{frontmatter}{%
\renewcommand{\headrulewidth}{0pt}%
\renewcommand{\footrulewidth}{0pt}%
\fancyhead{%
}%
\fancyfoot{%
}%
\fancyhf{}%
\fancyhead[L]{DV410}%
\fancyhead[C]{\thepage}%
\fancyhead[R]{23802}%
}%
\fancypagestyle{mainmatter}{%
\renewcommand{\headrulewidth}{0pt}%
\renewcommand{\footrulewidth}{0pt}%
\fancyhead{%
}%
\fancyfoot{%
}%
\fancyhf{}%
\fancyhead[L]{DV410}%
\fancyhead[C]{Page \thepage\ of \pageref{LastPage}}%
\fancyhead[R]{23802}%
}%
\pagenumbering{roman}%
\pagestyle{fancy}%
\fancyhf{}%
\fancyhead[L]{DV410}%
\fancyhead[C]{Page \thepage\ of \pageref{LastPage}}%
\fancyhead[R]{23802}%
%TC:ignore%
\title{
        \begin{flushright}
        \large \textbf{Candidate Number: 23802}
        \end{flushright}
        \vspace*{30mm}
        \begin{center}
        \large MSc in Development Management 2023 \\
        \vspace*{5mm}
        Dissertation submitted in partial fulfilment of the requirements of the degree. \\
        \vspace*{35mm}
        \Huge \textbf{Title title title} \\
        \vspace*{20mm}
        \end{center}
    }%
\date{}%
%TC:endignore%
%
\begin{document}%
\normalsize%
\includepdf[pages=-, offset=0 -1cm, frame]{DV410_Dissertation Cover Sheet_ Consent Form_Front Page_2023-24.pdf}%
\pagestyle{frontmatter}%
\maketitle%

\vfill
\begin{center}\textbf{Word Count: 3841}\end{center}
\newpage%
%TC:ignore%
\vspace*{\fill}%
\begin{center}%
\begin{minipage}{0.8\textwidth}%
\begin{center}%
\section*{Abstract}%
\end{center}%
\justify%
This is the abstract text. It should be centered on the page, and the text should be justified. This is the abstract text. It should be centered on the page, and the text should be justified. This is the abstract text. It should be centered on the page, and the text should be justified. This is the abstract text. It should be centered on the page, and the text should be justified. This is the abstract text. It should be centered on the page, and the text should be justified. %
\end{minipage}%
\end{center}%
\vspace*{\fill}%
%TC:endignore%
\newpage%
\tableofcontents%
\newpage%
%TC:ignore%
\section*{Abbreviations}%
\label{sec:Abbreviations}%
\begin{tabbing}%
\hspace{3cm} \= \kill%
\textbf{AI} \> Artificial Intelligence \\%
\textbf{API} \> Application Programming Interface \\%
\textbf{CPU} \> Central Processing Unit \\%
\textbf{GPU} \> Graphics Processing Unit \\%
\textbf{IoT} \> Internet of Things \\%
\textbf{ML} \> Machine Learning \\%
\textbf{NLP} \> Natural Language Processing \\%
\textbf{RAM} \> Random Access Memory \\%
\textbf{UI} \> User Interface \\%
\textbf{UX} \> User Experience \\%
\end{tabbing}

%
%TC:endignore%
\newpage%
%TC:ignore%
\listoffigures%
%TC:endignore%
\newpage%
%TC:ignore%
\listoftables%
%TC:endignore%
\newpage%
\pagenumbering{arabic}%
\pagestyle{mainmatter}%
\section{Introduction}%
\label{sec:Introduction}%
This is the introduction

The organization of this article is as follows: Section II provides a
brief literature review of the PTAs, South--South trade and the
importance of the structure of trade. Section III introduces the
methodology and data. Section IV presents the empirical results fol-
lowed by a discussion of the robustness tests. Section V concludes.
%
This is the introduction section. Here is a citation: %
\cite{dahi_preferential_2013}

%
\section{Literature Review}%
\label{sec:LiteratureReview}%
This section reviews the literature on the theoretical and empirical
effects of PTAs on exports and situates the analysis in the relevant
field of research.

\textbf{Theoretical Framework}

\textbf{Comparative advantage and trade creation and diversion}

Traditional trade theory emphasizes trade creation (allowing cheaper
products from PTA members to substitute for more expensive domestic
products) and trade diversion (substituting products from non-PTA
members that were cheaper before the PTA with products from PTA members
that are cheaper now due to the PTA reducing tariffs) (Schiff, Winters
and Schiff, 2003) and argues that the impact of PTAs depends on the
comparative advantage of member countries. In particular, it argues that
PTAs magnify the impacts of a country's comparative advantage, relative
to the world and to other member countries signatories of a common PTA.
If member countries of a PTA have a comparative advantage on a factor
endowment relative to the world, but one country also has a comparative
advantage on the same factor endowment relative to the other member
countries, the country with the ``extreme'' advantage will be more
vulnerable to trade diversion effects, while countries with
``intermediate'' advantages will gain from trade creation effects,
predicting divergence of trade outcomes, and winners and losers among
member countries. (Venables, 2003). This emphasis on the trade creation
and trade diversion effects among member countries with significant
differences in the comparative advantage of their factor endowments
relative to the world and to each other, suggests that, when the country
with the ``extreme'' comparative advantage is a high-income country,
relative to a lower-income country with an ``intermediate'' comparative
advantage, the lower-income country should seek a PTA with the other
country as it will gain more. On the contrary, if both members are
lower-income countries, the country with the ``extreme'' comparative
advantage, should not seek a PTA with the other low-income member
country as it will be vulnerable. (Sanguinetti, Siedschlag and
Martincus, 2010). This logic can be easily extended to the North-South
and South-South types of PTAs, as North countries will reasonably have
an ``extreme'' comparative advantage in skill-intensive goods relative
to South countries, while South countries will reasonably have an
``extreme'' comparative advantage in labour-intensive goods relative to
North countries. Furthermore, it is also argued that benefiting from
economies of scale through South-South economic integration is more
difficult because member countries do not have complementary production
and trade structures nor high interpenetration of each other's markets
on intra-industry trade. (Schiff, Winters and Schiff, 2003). Also, South
countries can benefit from greater technological diffusion from
North-South PTAs as the North countries have higher industrial
development as well as investment in research (Schiff and Wang, 2008).
Finally, as the trend in manufacturing has been in favour of vertical
specialization or value chain fragmentation (Krugman, 1995), North-South
PTAs are preferable as developing countries strive to capture a greater
portion of the value added. For these arguments, developing countries
should therefore be better off entering into North-South rather than
South-South agreements.

\textbf{Economies of Scale and Input-Output linkages}

In contrast, classical development theory and new trade literature go
beyond the static welfare gains from trade creation and diversion
effects when analysing the effect of PTAs.

``Preferential trade arrangements between developing countries can lead
to industrialization of the region as a whole as a consequence of the
effective market enlargement induced by reducing intra-South barriers.''
(Puga and Venables, 1998)

Other arguments

\begin{itemize}
\item
  Infant industry development
\item
  Economies of scale
\item
  Decoupling
\end{itemize}

``Trade can also arise from product differentiation and from economies
of scale that reduce costs as production grows. In these circumstances
competition between firms is weakened, and consumers lose. International
trade then offers an important means of increasing competition by
allowing new suppliers to enter markets.'' (Schiff, Winters and Schiff,
2003)

Stumbling block vs building block dichotomy

\textbf{Discussions on SS PTAs}

``\emph{The trade literature long argued that PTAs can benefit member
states through economies of scale and comparative advantage and higher
competition (Schiff, 2003). However, these arguments are generally
reserved for North--North and South--North but not South--South PTAs.
First, it is argued that similar production and trade structures in the
South make it more difficult to benefit from economies of scale. Second,
given the lower industrial development and research and development
activities in the South, greater technology diffusion for the Southern
country can be reaped from South--North integration (Schiff and Wang,
2008).3 Third, the more advanced members are argued to be the likely
winners in South--South integration, thanks to their higher industrial
and institutional development. As a result, lower income Southern
countries might be better off entering South-- North PTAs. It is also
claimed that industries with long term development potential are more
likely to move to the bigger and richer members, leading to divergence
once the barriers are lowered under South--South PTAs (Puga and
Venables, 1997; Schiff, 2003; Venables, 2003). Last but not the least,
North--South PTAs are argued to facilitate increasing vertical
specialization or value chain fragmentation, what Krugman (1995)
referred to as the slicing up of the value added.4}

\emph{In contrast, the classical development theory and new trade
literature has a more positive view of South--South PTAs, focusing on
their developmental benefits through infant industry development,
economies of scale and decoupling rather than on the static welfare
gains (from trade creation and diversion) or the `stumbling block/
building block' dichotomy. Myrdal (1956), for example, suggested that
regional integration in the South can help developing countries overcome
local market size limita- tions during industrialization. Accordingly,
given the strongly skill biased structure of output expansion in inter-
national trade (Antweiler and Trefler, 2002), increasing market size can
help developing countries enjoy scale effects and increase the skill
content of their exports while reducing the cost of intermediaries,
which in return may help stimulate export penetration into Northern mar-
kets in industrial goods (Fugazza and Robert-Nicoud, 2006). Likewise,
Lewis (1980), and more recently UNCTAD (2005) and World Bank (2008),
also pointed out that South--South trade can reduce the growth depen-
dence of the South on Northern growth, leading perhaps to decoupling
from Northern business cycles (thus helping the recovery from current
global downturn (ESCAP, 2009)). Furthermore, the structure of
South--South trade is argued to have dynamic and long term benefits for
developing countries because of its comparatively higher technology and
human capital intensive factor content (Amsden, 1987; Lall et al., 1989;
Demir and Dahi, 2011). Besides, similarity in production pattern and
resource base may facilitate appropriate technology trans- fer (Amsden,
1980, 1987; UNIDO, 2005; World Bank, 2006).}

\emph{It is also possible that South--North PTAs can yield more benefits
to Northern countries than the Southern ones because of asymmetries in
bargaining power, nego- tiating capacity and retaliatory capability.
Even though these asymmetries are also present between Southern
countries, the gap is likely to be smaller.}

\emph{Thrasher and Gallagher (2008), for example, show that South--South
PTAs leave the greatest policy space available to `deploy effective
policy for long-run diversification and develop- ment' than South--North
PTAs. We should also note that structuralist North--South models have
long discussed how interactions between countries with asymmetrical
economic structures, patterns of specialization, and devel- opment can
lead to uneven development (Findlay, 1980; Darity, 1990; Dutt, 1992; and
also see the survey articles Findlay, 1984; Dutt, 1989; Darity and
Davis, 2005).}

\emph{In addition to the debate above, the effects of PTAs on the
structure of trade are of particular importance for long term
development and growth. Development economics and the new trade theory
provide strong evidence that not all trade is equal and what you export
might matter for long term economic performance (Kaldor, 1967; An and
Iyigun, 2004; Hausmann et al., 2007). Exports in more technology
intensive industries are likely to generate lar- ger spillovers (such as
innovation and physical and human capital accumulation) and linkages for
development than lower technology and labour intensive ones (Hausman et
al., 2007). Earlier on, this point was also raised by Kaldor (1967) in
his three growth laws; which stated that there is a strong positive
relationship between the growth of manufacturing output and (i) the
growth of GDP, (ii) the growth of labour productivity in manufacturing
(i.e. the Verdoorn's law) and (iii) the growth of productivity in
nonmanufacturing sectors.}

\emph{Note that the question we raise here is different than the one
usually discussed in the literature, which is whether PTAs are trade
creating or diverting.5 To the extent that PTAs enhance manufactures
exports, we can then start evaluating the success or failure of PTAs
according to their potential long term developmental impact. Much of the
traditional PTA literature, both theoretical and empiri- cal, is taken
up by trade creation versus trade diversion debate. These questions are
not unimportant; however, there is reason to question the
disproportionate attention still given to the classic Vinerian
dichotomy.}

\emph{(...)}

\emph{Second, since North--North, South--North and North--South trade
barriers appear to be significantly lower than the ones present in
South--South trade (Kowalski and Shepherd, 2006, also see Kee et al.,
2009; Medvedev, 2010), it is unlikely that South--South PTAs are trade
diverting from the North, which has retrospectively been the main point
of contention among trade theorists on the relative costs and benefits
of South--South PTAs. In fact, consistent with Mundell (1968)'s
assertion that `a member's gain from a free-trade area will be larger,
the higher are the initial tariffs of partner countries', South--South
trade barrier reduction is found to generate a significant increase in
South--South exports, while no such effect is reported in the case of
North--South, South--North or North--North trade (Kowalski and Shepherd,
2006). Besides, there is also some empirical evidence showing that
South--South PTAs are no more trade diverting than other PTAs (Cernat,
2001). Third, since higher transportation costs and former colonial
linkages with Northern countries (which always appear to be significant
in gravity models of trade), in addition to higher trade barriers (Kee
et al., 2009), con- tinue to limit South--South trade expansion, PTAs
might be seen as a way of compensating for such trade barriers that are
lower in South--North, North--South or North-- North trade.6 Last but
not least, in the case of industrial development, what matters are
dynamic not static gains. That is to say, if South--South PTAs are found
to enhance industrial development, the long term gains may very well
outweigh the static short term losses.}'' (Dahi and Demir, 2013)

\textbf{Significance of Exports}

\textbf{Defining South and North}

\textbf{Empirical literature}

``Turning to the empirical work on PTAs, the majority of research
reports a significantly positive effect of PTAs on member trade.
Cipollina and Salvatici (2010) review 85 papers on the effects of PTAs
and find that the mean effect is 0.59 (or an 80\% increase in trade),
while the median is 0.38 (or a 46\% increase in trade). Although the
range of coefficient estimates is quite large (--9.01 to 15.41), only
312 out of 1827 coefficient estimates are reported as negative.
Nevertheless, despite the diversity of research, there are only few
studies that compare heterogeneous effects of PTAs within and between
developing and developed countries. Among the few, Medvedev (2010),
using a cross sectional analysis, reports that while North--North PTAs
are insignificant in stimulating preferential trade, North--South PTAs
increase trade by 40\% and South-- South PTAs increase them by 163\%.

Cipollina, M. and Salvatici, L. (2010) Reciprocal trade agreements in
gravity models: a meta analysis, Review of International Economics, 18,
63-80.

Medvedev, Denis, Preferential Trade Agreements and Their Role in World
Trade (October 1, 2006). World Bank Policy Research Working Paper No.
4038, Available at SSRN: \url{https://ssrn.com/abstract=938031}''

The empirical work on the structure of trade under PTAs has also been
scarce. Sanguinetti et al. (2010) examine the impact of PTAs on
South--South manufacturing production patterns in the case of MERCOSUR
for the period of 1985--1998 and find that South--South PTAs cause a
spatial regional reorganization of production along the lines of
internal comparative advantage.

Sanguinetti, P., Siedschlag, I., \& Martincus, C. V. (2010). The Impact
of South-South Preferential Trade Agreements on Industrial Development:
An Empirical Test. Journal of Economic Integration, 25(1), 69--103.
http://www.jstor.org/stable/23000966

\textbf{Stylised facts}

\textbf{Ideas:}

Total number of SS, NS and NN PTAs

Share of SS, NS and NN PTAs

Total exports by S and N countries

Share of total exports by S and N countries

Total exports of manufactured by S and N countries

Share of total exports of manufactured by S and N countries

Number of products exported by S and N countries

\textbf{References}

Dahi, O.S. and Demir, F. (2013) `Preferential trade agreements and
manufactured goods exports: does it matter whom you PTA with?',
\emph{Applied Economics}, 45(34), pp. 4754--4772. Available at:
https://doi.org/10.1080/00036846.2013.804169.

Krugman, P. (1995) `Growing World Trade: Causes and Consequences',
\emph{Brookings Papers on Economic Activity} {[}Preprint{]}.

Puga, D. and Venables, A.J. (1998) `Trading Arrangements and Industrial
Development'.

Sanguinetti, P., Siedschlag, I. and Martincus, C.V. (2010) `The Impact
of South-South Preferential Trade Agreements on Industrial Development:
An Empirical Test', \emph{Journal of Economic Integration}, 25(1), pp.
69--103.

Schiff, M. and Wang, Y. (2008) `North-South and South-South
Trade-Related Technology Diffusion: How Important Are They in Improving
TFP Growth?', \emph{The Journal of Development Studies}, 44(1), pp.
49--59. Available at: https://doi.org/10.1080/00220380701722282.

Schiff, M.W., Winters, L.A. and Schiff, M. (2003) \emph{Regional
Integration And Development}. Washington, UNITED STATES: World Bank
Publications. Available at:
http://ebookcentral.proquest.com/lib/londonschoolecons/detail.action?docID=3050563
(Accessed: 12 August 2024).

Venables, A.J. (2003) `Winners and Losers from Regional Integration
Agreements', \emph{The Economic Journal}, 113(490), pp. 747--761.
Available at: https://doi.org/10.1111/1468-0297.t01-1-00155.

%
\section{Methodology}%
\label{sec:Methodology}%
The Gravity Model of Trade

Often referred as the ``workhorse'' of international trade, the gravity
model is prominent in the empirical literature of applied international
trade analysis. Among the arguments that could explain/support the use
of the gravity model, there are four that are particularly relevant for
our purposes. First, the gravity model of trade is intuitive to
understand. Following the metaphor of Newton's Law of Universal
Gravitation, it predicts that international trade between two countries
is directly proportional to the product of their economic size, and
inversely proportional to trade frictions between them. In simpler
words, the bigger (smaller) the economies of two countries, and the
easier (harder) it is for them to trade with each other, the more (less)
we expect them to trade. Second, it is referred to as a structural model
with solid theoretical foundations, which makes it appropriate for
counterfactual analysis, such as measuring the effects of trade policies
as we aim to do with the effects of South-North versus South-South
agreements. Third, model has a flexible structure, which will allow us
to construct a specification tailored to our research. Finally, fourth,
it holds consistent and remarkable predictive power, both with aggregate
and sectoral data (Yotov et al. 2016).

Through the decades, the gravity equation has been regularly upgraded in
the theoretical and empirical literature. Of relevance, the simple
intuition of the gravity model was theoretically extended by Anderson to
note that, after controlling for size, the increase or decrease is
\emph{relative} to the average barriers of the two countries with all
their partners, which are referred as ``multilateral resistance''
(Anderson 1979). The more trade barriers or resistance to trade exists
with other countries relative to a given partner, the more a country is
pushed to trade with said partner. Anderson also introduced the
assumptions of product differentiation by place of origin, and Constant
Elasticity of Substitution (CES) expenditures ,or the Armington-CES
assumption (Yotov et al. 2016; Chatzilazarou and Dadakas 2023), which
led us to today's generalized form of the gravity equation, as developed
and popularised by Anderson and van Wincoop (Anderson and van Wincoop
2003).

Do I need to add more on the theory underlying the Gravity Models of
Trade?

Equally important, several empirical developments have strengthened the
gravity model and inform our choice of methodology: Exporter-time and
importer-time fixed effects are used to account for the multilateral
resistance terms in a gravity estimation with panel data (Olivero and
Yotov 2012); As the gravity model is often estimated with an OSL
estimator, zero-trade flows were dropped from the sample when trade was
transformed into a logarithmic form. Also, trade data is recognized to
suffer from heteroscedasticity (Yotov et al. 2016). To solve for
zero-trade flows and heteroscedasticity, the Poisson Pseudo Maximum
Likelihood (PPML) estimator has been proposed to estimate the gravity
model, avoiding potential biases (Silva and Tenreyro 2006; Santos Silva
and Tenreyro 2011); Country-pair fixed effects has been proposed to
account for the unobserved endogeneity of trade policy (Baier and
Bergstrand 2007). It is worth nothing that the inclusion of
exporter-time and importer-time fixed effects will absorb all observable
and unobservable time-varying country-specific characteristics that
could affect the dependent variable, while the country-pair fixed
effects will absorb observable and unobservable bilateral time-invariant
characteristics that could affect trade costs; The inclusion of
intra-trade flows as well as international trade flows is proposed to
correctly estimate the effects of non-discriminatory trade policy,
allowing for consumers to choose products from both international and
domestic sources (Dai, Yotov, and Zylkin 2014; Heid, Larch, and Yotov
2017); Year-intervals instead of data pooled over consecutive years
should be used to allow for adjustment of trade flows to policies that
might not have immediate effects (Baier and Bergstrand 2007; Anderson
and Yotov 2016); And finally, to account for the effects of
globalization forces that may biased the estimates of trade policies, a
set of globalization dummies are recommended to control for the effects
of globalization in the gravity model (Yotov 2012; Bergstrand, Larch,
and Yotov 2015). Based on the theoretical and empirical best-practices
found in the relevant literature, we employ the following gravity
equation using a PPML estimator and a balanced panel data approach with
multiple exporters, multiple importers and time as our benchmark model:

\[(1)\ \ \ \ \ X_{ij,t} = \ exp(\eta_{i,t} + \psi_{j,t} + \gamma_{\binom{-}{ij}} + \beta_{1}{PTA}_{ij,t} + \beta_{2}{PTA}_{ij,t - 5} + \sum_{t}^{}b_{t}) + \epsilon_{ij,t}\]

Where \(X_{ij,t}\) denotes the value of exports from an origin country
\(i\) to a destination country \(j\); \(\eta_{i,t}\) and \(\psi_{j,t}\)
are, respectively, exporter-time and importer-time fixed-effects;
\(\gamma_{\binom{-}{ij}}\) is a country-pair fixed-effect;
\({PTA}_{ij,t}\) and \({PTA}_{ij,t - 5}\) are our main variables of
interest, which, respectively indicate if \(i\) and \(j\) are members of
a PTA at time \(t\) and, to account for potential ``phase-in'' effects
over time of the PTA, at time \(t - 5\); \(\sum_{t}^{}b_{t}\) is a set
of dummies that equal 1 for international trade and 0 for domestic trade
observations at each time \(t\); and \(\epsilon_{ij,t}\) is an error
term.

In contrast with our main interest of research, which are the potential
heterogenous effects of PTAs on different members for different types of
agreements, this benchmark model, specifically
\(\beta = \beta_{1} + \beta_{2}\) , would provide the average ``total''
partial effect of PTAs on trade after accounting for lagged effects, but
it cannot provide the effects for a given agreement, country-pairs o
specific country members to a specific agreement. As such, three
successive expansions can be implemented to capture heterogeneity in PTA
effects as proposed by Baier \emph{et al}. (Baier, Yotov, and Zylkin
2019):

\[(2)\ \ \ \ \ X_{ij,t} = \ exp(\eta_{i,t} + \psi_{j,t} + \gamma_{\binom{-}{ij}} + \sum_{A}^{}{\beta_{1,A}{PTA}_{ij,t}} + \sum_{A}^{}{\beta_{2,A}{PTA}_{ij,t - 5}} + \sum_{t}^{}b_{t}) + \epsilon_{ij,t}\]

Equation (2) can be implemented to account for heterogeneous effects of
PTAs at the level of the specific agreement, by allowing for distinct
average partial effects for each individual agreement, using superscript
\(A\) to index by agreement and also allowing for agreement-specific
lags: \(\beta_{A} = \beta_{1,A} + \beta_{2,A}\).

Extended Benchmark Model with PTA Types

You want to modify the benchmark model to allow for different types of
PTAs (NN, NS, SS) while still keeping the same overall structure. Here's
how you can adjust the model:

Equation (1) with PTA Types:

\[(3)\ \ \ \ \ X_{ij,t} = \ exp(\eta_{i,t} + \psi_{j,t} + \gamma_{\binom{-}{ij}} + \beta_{1NN}{PTA\_ NN}_{ij,t} + \beta_{2NN}{PTA\_ NN}_{ij,t - 5} + + \beta_{1NS}{PTA\_ NS}_{ij,t} + \beta_{2NS}{PTA\_ NS}_{ij,t - 5} + \beta_{1SS}{PTA\_ SS}_{ij,t} + \beta_{2SS}{PTA\_ SS}_{ij,t - 5} + \sum_{t}^{}b_{t}) + \epsilon_{ij,t}\]

Where:

\begin{itemize}
\item
  \(X_{ij,t}\)\hspace{0pt} denotes the value of exports from country
  \(i\) to country \(j\) at time \(t\).
\item
  \(\eta_{i,t}\) and \(\psi_{j,t}\ \)are exporter-time and importer-time
  fixed effects, respectively.
\item
  \(\gamma_{\binom{-}{ij}}\) is a country-pair fixed effect.
\item
  \hspace{0pt}\(\beta_{1NN}\) and \(\beta_{2NN}\) are the coefficients
  for the immediate and lagged effects of a North-North PTA
  (\(PTA\_ NN\)).
\item
  \hspace{0pt}\hspace{0pt}\(\beta_{1NS}\) and \(\beta_{2NS}\) are the
  coefficients for the immediate and lagged effects of a North-South PTA
  (\(PTA\_ SN\)).
\item
  \hspace{0pt}\hspace{0pt}\(\beta_{1SS}\) and \(\beta_{2SS}\) are the
  coefficients for the immediate and lagged effects of a South-South PTA
  (\(PTA\_ SS\)).
\item
  \(\sum_{t}^{}b_{t}\) is a set of time dummies accounting for
  international trade-specific effects at each time \(t\).
\item
  \(\epsilon_{ij,t}\) is the error term.
\end{itemize}

Extended Model with PTA Heterogeneity and Types

Next, you want to extend the model allowing for PTA heterogeneity to
also distinguish between different types of PTAs.

Equation (2) with PTA Types:

\[(4)\ \ \ \ \ X_{ij,t} = \ exp(\eta_{i,t} + \psi_{j,t} + \gamma_{\binom{-}{ij}} + \sum_{A}^{}{(\beta_{1,A,NN}{PTA\_ NN}_{ij,t}\  + \ \beta_{2,A,NN}{PTA\_ NN}_{ij,t - 5})} + \sum_{A}^{}{(\beta_{1,A,NS}{PTA\_ NS}_{ij,t}\  + \ \beta_{2,A,NS}{PTA\_ NS}_{ij,t - 5})} + \sum_{A}^{}{(\beta_{1,A,SS}{PTA\_ SS}_{ij,t}\  + \ \beta_{2,A,SS}{PTA\_ SS}_{ij,t - 5})} + \sum_{t}^{}b_{t}) + \epsilon_{ij,t}\]

Where:

\begin{itemize}
\item
  \(X_{ij,t}\)\hspace{0pt} denotes the value of exports from country
  \(i\) to country \(j\) at time \(t\).
\item
  \(\eta_{i,t}\) and \(\psi_{j,t}\ \)are exporter-time and importer-time
  fixed effects, respectively.
\item
  \(\gamma_{\binom{-}{ij}}\) is a country-pair fixed effect.
\item
  The summations \hspace{0pt}\(\sum_{}^{}A\) denote the sum over
  different agreements \(A\) for:

  \begin{itemize}
  \item
    \(\beta_{1,A,NN}\) and \(\beta_{2,A,NN}\): Coefficients for the
    immediate and lagged effects of North-North PTAs
    \hspace{0pt}(\(PTA\_ NN\)).
  \item
    \(\beta_{1,A,NS}\) and \(\beta_{2,A,NS}\): Coefficients for the
    immediate and lagged effects of North-South PTAs (\(PTA\_ SN\)).
  \item
    \(\beta_{1,A,SS}\) and \(\beta_{2,A,SS}\): Coefficients for the
    immediate and lagged effects of South-South PTAs (\(PTA\_ SS\)).
  \end{itemize}
\item
  \(\sum_{t}^{}b_{t}\) is a set of time dummies accounting for
  trade-specific effects at each time \(t\).
\item
  \(\epsilon_{ij,t}\) is the error term.
\end{itemize}

Variables:

\begin{enumerate}
\def\labelenumi{\arabic{enumi}.}
\item
  \({PTA\_ NN}_{ij,t}\): Dummy variable that takes the value of 1 if the
  trade pair \((i,j)\) is part of a North-North PTA at time \(t\), and 0
  otherwise.
\item
  \({PTA\_ NN}_{ij,t - 5}\): Dummy variable that takes the value of 1 if
  the trade pair \((i,j)\) was part of a North-North PTA at time
  \(t\)\emph{-5}, and 0 otherwise.
\item
  \({PTA\_ NS}_{ij,t}\): Dummy variable that takes the value of 1 if the
  trade pair \((i,j)\) is part of a North-South PTA at time \(t\), and 0
  otherwise.
\item
  \({PTA\_ NS}_{ij,t - 5}\): Dummy variable that takes the value of 1 if
  the trade pair \((i,j)\) was part of a North-South PTA at time
  \(t\)\emph{-5}, and 0 otherwise.
\item
  \({PTA\_ SS}_{ij,t}\): Dummy variable that takes the value of 1 if the
  trade pair \((i,j)\) is part of a South-South PTA at time \(t\), and 0
  otherwise.
\item
  \({PTA\_ SS}_{ij,t - 5}\):: Dummy variable that takes the value of 1
  if the trade pair \((i,j)\) was part of a South-South PTA at time
  \(t\)\emph{-5}, and 0 otherwise.
\end{enumerate}

These models allow you to capture the differentiated impacts of PTAs
depending on whether they are between two developed countries (NN),
between a developed and a developing country (NS), or between two
developing countries (SS).

References

%
\section{Findings}%
\label{sec:Findings}%
These are the findings

%
\subsection{Benchmark Estimation Results by Region}%
\label{subsec:BenchmarkEstimationResultsbyRegion}%
\subsubsection{Benchmark Model Results by Region}%
\label{ssubsec:BenchmarkModelResultsbyRegion}%
\begin{table}[htbp]
    \centering
    \caption{Benchmark Model Regional Results}
    \label{tab:benchmark_region_analysis} % This allows you to reference the table in the text with \ref{tab:pta_analysis}
    \begin{adjustbox}{max width=\textwidth}
    \begin{tabular}{l@{\extracolsep{1pt}}ccccc}
    \hline
    & \multicolumn{1}{c}{(1)} & \multicolumn{1}{c}{(2)} & \multicolumn{1}{c}{(3)} & \multicolumn{1}{c}{(4)} & \multicolumn{1}{c}{(5)} \\
    \hline
    \textbf{Variables} &  &  &  &  &  \\
    \hline
     & PPML & PPML & PPML & PPML & PPML \\
     & Africa & Americas & Asia & Europe & Intercontinental \\
    \hline
    PTA & 0.578$^{\ast\ast\ast}$ & 0.287$^{\ast\ast\ast}$ & 0.064 & 0.237$^{\ast\ast\ast}$ & 0.015 \\
    & (0.154) & (0.071) & (0.083) & (0.019) & (0.093) \\

    PTA Lag & -0.278 & 0.146 & -0.167$^{\ast\ast\ast}$ & 0.238$^{\ast\ast\ast}$ & 0.188$^{\ast\ast\ast}$ \\
    & (0.300) & (0.149) & (0.056) & (0.022) & (0.043) \\

    PTA + PTA Lag & 0.301 & 0.433$^{\ast\ast\ast}$ & -0.103 & 0.475$^{\ast\ast\ast}$ & 0.203$^{\ast}$ \\
    & (0.295) & (0.140) & (0.094) & (0.025) & (0.106) \\
    \hline
    Exporter-Year FE & Yes & Yes & Yes & Yes & Yes \\
    Importer-Year FE & Yes & Yes & Yes & Yes & Yes \\
    Country-Pair FE & Yes & Yes & Yes & Yes & Yes \\
    R-Squared & 0.997 & 0.999 & 0.999 & 0.997 & 0.998 \\
    Observations & 5838 & 10997 & 25308 & 28168 & 73930 \\
    \hline
    \multicolumn{6}{l}{\footnotesize{Notes: Robust standard errors clustered at the country-pair in parentheses. Significance levels are indicated as follows: $^{\ast}$p$<$0.1; $^{\ast\ast}$p$<$0.05; $^{\ast\ast\ast}$p$<$0.01.}} \\
    \end{tabular}
    \end{adjustbox}
\end{table}    

%
\subsection{PTA Estimation Results by Region}%
\label{subsec:PTAEstimationResultsbyRegion}%
\subsubsection{PTA Model Results by Region}%
\label{ssubsec:PTAModelResultsbyRegion}%
\begin{table}[htbp]
    \centering
    \caption{PTA + PTA Lag Coefficients for Africa Region}
    \label{tab:pta_africa}
    \begin{adjustbox}{max width=\textwidth}
    \begin{tabular}{lcc}
    \hline
    \textbf{Statistically Insignificant} &  &  \\
    \hline
    PTA ID & Estimate & SE \\
    \hline
    670 & 0.326 & (0.410) \\
    787 & 0.304 & (0.233) \\
    \hline
    \multicolumn{3}{l}{\footnotesize{Notes: Robust standard errors clustered at the country-pair level in parentheses.}} \\
    \multicolumn{3}{l}{\footnotesize{Significance levels are indicated as follows: $^{\ast}$p$<$0.1; $^{\ast\ast}$p$<$0.05; $^{\ast\ast\ast}$p$<$0.01.}} \\
    \end{tabular}
    \end{adjustbox}
\end{table}
%
\begin{table}[htbp]
    \centering
    \caption{PTA + PTA Lag Coefficients for Americas Region}
    \label{tab:pta_americas}
    \begin{adjustbox}{max width=\textwidth}
    \begin{tabular}{lcc}
    \hline
    \textbf{Positive and Statistically Significant} &  &  \\
    \hline
    PTA ID & Estimate & SE \\
    \hline
    213 & 1.342$^{\ast\ast\ast}$ & (0.434) \\
    218 & 0.879$^{\ast\ast\ast}$ & (0.173) \\
    716 & 0.732$^{\ast\ast}$ & (0.358) \\
    239 & 0.571$^{\ast\ast\ast}$ & (0.173) \\
    201 & 0.545$^{\ast\ast}$ & (0.265) \\
    612 & 0.515$^{\ast\ast}$ & (0.251) \\
    616 & 0.488$^{\ast\ast\ast}$ & (0.044) \\
    168 & 0.410$^{\ast\ast\ast}$ & (0.113) \\
    163 & 0.342$^{\ast\ast\ast}$ & (0.096) \\
    141 & 0.265$^{\ast\ast\ast}$ & (0.024) \\
    \hline
    \textbf{Statistically Insignificant} &  &  \\
    \hline
    PTA ID & Estimate & SE \\
    \hline
    185 & 0.291 & (0.376) \\
    645 & 0.117 & (0.141) \\
    \hline
    \textbf{Negative and Statistically Significant} &  &  \\
    \hline
    PTA ID & Estimate & SE \\
    \hline
    188 & -0.774$^{\ast\ast\ast}$ & (0.144) \\
    \hline
    \multicolumn{3}{l}{\footnotesize{Notes: Robust standard errors clustered at the country-pair level in parentheses.}} \\
    \multicolumn{3}{l}{\footnotesize{Significance levels are indicated as follows: $^{\ast}$p$<$0.1; $^{\ast\ast}$p$<$0.05; $^{\ast\ast\ast}$p$<$0.01.}} \\
    \end{tabular}
    \end{adjustbox}
\end{table}
%
\begin{table}[htbp]
    \centering
    \caption{PTA + PTA Lag Coefficients for Asia Region}
    \label{tab:pta_asia}
    \begin{adjustbox}{max width=\textwidth}
    \begin{tabular}{lcc}
    \hline
    \textbf{Positive and Statistically Significant} &  &  \\
    \hline
    PTA ID & Estimate & SE \\
    \hline
    683 & 1.080$^{\ast\ast\ast}$ & (0.237) \\
    70  & 0.472$^{\ast\ast\ast}$ & (0.150) \\
    100 & 0.376$^{\ast\ast\ast}$ & (0.105) \\
    67  & 0.342$^{\ast\ast\ast}$ & (0.125) \\
    675 & 1.360$^{\ast\ast}$ & (0.655) \\
    475 & 0.636$^{\ast\ast}$ & (0.298) \\
    598 & 0.166$^{\ast\ast}$ & (0.083) \\
    474 & 0.419$^{\ast}$ & (0.243) \\
    \hline
    \textbf{Statistically Insignificant} &  &  \\
    \hline
    PTA ID & Estimate & SE \\
    \hline
    72  & 0.254 & (0.178) \\
    116 & 0.256 & (0.703) \\
    492 & 0.041 & (0.180) \\
    640 & 0.183 & (0.217) \\
    223 & -0.014 & (0.203) \\
    71  & -0.138 & (0.091) \\
    456 & -0.209 & (0.165) \\
    534 & -0.165 & (0.370) \\
    667 & -0.049 & (0.241) \\
    \hline
    \textbf{Negative and Statistically Significant} &  &  \\
    \hline
    PTA ID & Estimate & SE \\
    \hline
    221 & -2.955$^{\ast\ast\ast}$ & (0.727) \\
    220 & -1.215$^{\ast\ast\ast}$ & (0.093) \\
    599 & -0.967$^{\ast\ast\ast}$ & (0.191) \\
    1   & -0.732$^{\ast\ast}$ & (0.359) \\
    \hline
    Exporter-Year FE & Yes \\
    Importer-Year FE & Yes \\
    Country-Pair FE & Yes \\
    R-Squared & 0.999 \\
    Observations & 25308 \\
    \hline
    \multicolumn{3}{l}{\footnotesize{Notes: Robust standard errors clustered at the country-pair level in parentheses.}} \\
    \multicolumn{3}{l}{\footnotesize{Significance levels are indicated as follows: $^{\ast}$p$<$0.1; $^{\ast\ast}$p$<$0.05; $^{\ast\ast\ast}$p$<$0.01.}} \\
    \end{tabular}
    \end{adjustbox}
\end{table}
%
\begin{table}[htbp]
    \centering
    \caption{PTA + PTA Lag Coefficients for Europe Region}
    \label{tab:pta_europe}
    \begin{adjustbox}{max width=\textwidth}
    \begin{tabular}{lcc}
    \hline
    \textbf{Positive and Statistically Significant} &  &  \\
    \hline
    PTA ID & Estimate & SE \\
    \hline
    5   & 3.812$^{\ast\ast\ast}$ & (0.278) \\
    128 & 2.712$^{\ast\ast\ast}$ & (0.211) \\
    13  & 2.256$^{\ast\ast\ast}$ & (0.262) \\
    132 & 2.241$^{\ast\ast\ast}$ & (0.252) \\
    192 & 1.107$^{\ast\ast\ast}$ & (0.163) \\
    7   & 1.153$^{\ast\ast\ast}$ & (0.272) \\
    328 & 0.671$^{\ast\ast\ast}$ & (0.175) \\
    8   & 0.667$^{\ast\ast\ast}$ & (0.161) \\
    621 & 0.618$^{\ast\ast\ast}$ & (0.186) \\
    135 & 0.615$^{\ast\ast\ast}$ & (0.217) \\
    254 & 0.565$^{\ast\ast\ast}$ & (0.084) \\
    394 & 0.745$^{\ast\ast\ast}$ & (0.202) \\
    335 & 0.472$^{\ast\ast\ast}$ & (0.025) \\
    129 & 0.553$^{\ast\ast\ast}$ & (0.206) \\
    9   & 0.580$^{\ast\ast}$ & (0.285) \\
    11  & 0.656$^{\ast\ast}$ & (0.307) \\
    131 & 0.615$^{\ast\ast}$ & (0.281) \\
    594 & 0.474$^{\ast}$ & (0.251) \\
    \hline
    \textbf{Statistically Insignificant} &  &  \\
    \hline
    PTA ID & Estimate & SE \\
    \hline
    6   & 0.355 & (0.358) \\
    150 & 0.247 & (0.687) \\
    153 & 0.614 & (0.633) \\
    154 & 0.592 & (0.409) \\
    255 & 0.167 & (0.237) \\
    389 & 0.412 & (0.323) \\
    331 & 0.142 & (0.201) \\
    12  & -0.246 & (1.208) \\
    156 & -0.441 & (0.445) \\
    \hline
    \textbf{Negative and Statistically Significant} &  &  \\
    \hline
    PTA ID & Estimate & SE \\
    \hline
    133 & -0.772$^{\ast\ast\ast}$ & (0.248) \\
    \hline
    Exporter-Year FE & Yes \\
    Importer-Year FE & Yes \\
    Country-Pair FE & Yes \\
    R-Squared & 0.997 \\
    Observations & 28168 \\
    \hline
    \multicolumn{3}{l}{\footnotesize{Notes: Robust standard errors clustered at the country-pair level in parentheses.}} \\
    \multicolumn{3}{l}{\footnotesize{Significance levels are indicated as follows: $^{\ast}$p$<$0.1; $^{\ast\ast}$p$<$0.05; $^{\ast\ast\ast}$p$<$0.01.}} \\
    \end{tabular}
    \end{adjustbox} 
\end{table}
%
\begin{center}
\small
\begin{longtable}{lcc}
    \caption{PTA+PTALag Coefficients for Intercontinental} \label{tab:pta_intercontinental} \\
    
    \hline
    \textbf{Positive and Statistically Significant} &  &  \\
    \hline
    PTA ID & Estimate & SE \\
    \hline
    \endfirsthead
    
    \multicolumn{3}{c}{{\bfseries \tablename\ \thetable{} -- continued from previous page}} \\
    \hline
    PTA ID & Estimate & SE \\
    \hline
    \endhead
    
    \hline \multicolumn{3}{r}{{Continued on next page}} \\ \hline
    \endfoot
    
    \hline
    \endlastfoot
    
    627 & 2.372$^{\ast\ast\ast}$ & (0.345) \\
    415 & 1.853$^{\ast\ast\ast}$ & (0.201) \\
    206 & 1.539$^{\ast\ast\ast}$ & (0.180) \\
    75  & 1.366$^{\ast\ast\ast}$ & (0.493) \\
    263 & 1.426$^{\ast\ast\ast}$ & (0.115) \\
    4   & 1.254$^{\ast\ast\ast}$ & (0.268) \\
    626 & 1.099$^{\ast\ast\ast}$ & (0.121) \\
    657 & 0.705$^{\ast\ast\ast}$ & (0.082) \\
    637 & 0.667$^{\ast\ast\ast}$ & (0.102) \\
    202 & 0.658$^{\ast\ast\ast}$ & (0.123) \\
    208 & 0.763$^{\ast\ast\ast}$ & (0.129) \\
    136 & 0.744$^{\ast\ast\ast}$ & (0.185) \\
    490 & 0.843$^{\ast\ast\ast}$ & (0.181) \\
    17  & 0.811$^{\ast\ast\ast}$ & (0.242) \\
    466 & 0.710$^{\ast\ast\ast}$ & (0.147) \\
    304 & 0.770$^{\ast\ast\ast}$ & (0.120) \\
    628 & 0.484$^{\ast\ast\ast}$ & (0.142) \\
    207 & 0.516$^{\ast\ast\ast}$ & (0.114) \\
    518 & 0.627$^{\ast\ast\ast}$ & (0.135) \\
    330 & 0.314$^{\ast\ast\ast}$ & (0.086) \\
    164 & 0.288$^{\ast\ast\ast}$ & (0.073) \\
    96 & 0.271$^{\ast\ast\ast}$ & (0.055) \\
    181 & 0.392$^{\ast\ast}$ & (0.178) \\
    624 & 0.388$^{\ast\ast}$ & (0.163) \\
    521 & 0.101$^{\ast\ast}$ & (0.045) \\
    384  & 0.645$^{\ast}$ & (0.355) \\
    15  & 0.313$^{\ast}$ & (0.179) \\
    227 & 0.348$^{\ast}$ & (0.186) \\
    \hline
    \textbf{Statistically Insignificant} &  &  \\
    \hline
    PTA ID & Estimate & SE \\
    \hline
    641 & 2.028 & (1.255) \\
    543 & 1.090 & (0.707) \\
    509 & 0.210 & (0.216) \\
    252 & 0.192 & (0.357) \\
    508 & 0.140 & (0.122) \\
    376 & 0.172 & (0.228) \\
    416 & 0.424 & (0.295) \\
    401 & 0.407 & (0.288) \\
    152 & 0.110 & (0.266) \\
    242 & 0.050 & (0.294) \\
    390 & 0.0471 & (0.181) \\
    396 & 0.019 & (0.379) \\
    205 & 0.0012 & (0.178) \\
    602 & -0.076 & (0.918) \\
    383 & -0.202 & (0.152) \\
    386 & -0.092 & (0.168) \\
    84  & -0.059 & (0.120) \\
    979  & -0.126 & (0.294) \\
    644  & -0.189 & (0.122) \\
    658  & -0.303 & (0.349) \\
    \hline
    \textbf{Negative and Statistically Significant} &  &  \\
    \hline
    PTA ID & Estimate & SE \\
    \hline
    399 & -0.473$^{\ast\ast\ast}$ & (0.127) \\
    104 & -0.338$^{\ast\ast\ast}$ & (0.112) \\
    677 & -1.366$^{\ast\ast\ast}$ & (0.385) \\
    679 & -1.429$^{\ast\ast\ast}$ & (0.430) \\
    323 & -0.338$^{\ast\ast}$ & (0.138) \\
    512 & -0.458$^{\ast}$ & (0.266) \\
    \hline
    Exporter-Year FE & Yes \\
    Importer-Year FE & Yes \\
    Country-Pair FE & Yes \\
    R-Squared & 0.998 \\
    Observations & 73930 \\
    \hline
    \multicolumn{3}{l}{\footnotesize{Notes: Robust standard errors clustered at the country-pair level in parentheses.}} \\
    \multicolumn{3}{l}{\footnotesize{Significance levels are indicated as follows: $^{\ast}$p$<$0.1; $^{\ast\ast}$p$<$0.05; $^{\ast\ast\ast}$p$<$0.01.}} \\
\end{longtable}
\end{center}


% \begin{table}[htbp]
%     \centering
%     \caption{PTA + PTA Lag Coefficients for Intercontinental Region}
%     \label{tab:pta_intercontinental}
%     \begin{minipage}[t]{0.48\textwidth} % Adjust the width as needed
%         \centering
%         \begin{adjustbox}{max width=\textwidth}
%         \begin{tabular}{lcc}
%         \hline
%         \textbf{Positive and Statistically Significant} &  &  \\
%         \hline
%         PTA ID & Estimate & SE \\
%         \hline
%         627 & 2.372$^{\ast\ast\ast}$ & (0.345) \\
%         415 & 1.853$^{\ast\ast\ast}$ & (0.201) \\
%         206 & 1.539$^{\ast\ast\ast}$ & (0.180) \\
%         75  & 1.366$^{\ast\ast\ast}$ & (0.493) \\
%         263 & 1.426$^{\ast\ast\ast}$ & (0.115) \\
%         4   & 1.254$^{\ast\ast\ast}$ & (0.268) \\
%         626 & 1.099$^{\ast\ast\ast}$ & (0.121) \\
%         657 & 0.705$^{\ast\ast\ast}$ & (0.082) \\
%         637 & 0.667$^{\ast\ast\ast}$ & (0.102) \\
%         202 & 0.658$^{\ast\ast\ast}$ & (0.123) \\
%         208 & 0.763$^{\ast\ast\ast}$ & (0.129) \\
%         136 & 0.744$^{\ast\ast\ast}$ & (0.185) \\
%         490 & 0.843$^{\ast\ast\ast}$ & (0.181) \\
%         17  & 0.811$^{\ast\ast\ast}$ & (0.242) \\
%         466 & 0.710$^{\ast\ast\ast}$ & (0.147) \\
%         304 & 0.770$^{\ast\ast\ast}$ & (0.120) \\
%         628 & 0.484$^{\ast\ast\ast}$ & (0.142) \\
%         207 & 0.516$^{\ast\ast\ast}$ & (0.114) \\
%         518 & 0.627$^{\ast\ast\ast}$ & (0.135) \\
%         330 & 0.314$^{\ast\ast\ast}$ & (0.086) \\
%         164 & 0.288$^{\ast\ast\ast}$ & (0.073) \\
%         96 & 0.271$^{\ast\ast\ast}$ & (0.055) \\
%         181 & 0.392$^{\ast\ast}$ & (0.178) \\
%         624 & 0.388$^{\ast\ast}$ & (0.163) \\
%         521 & 0.101$^{\ast\ast}$ & (0.045) \\
%         384  & 0.645$^{\ast}$ & (0.355) \\
%         15  & 0.313$^{\ast}$ & (0.179) \\
%         227 & 0.348$^{\ast}$ & (0.186) \\
%         \hline
%         \end{tabular}
%         \end{adjustbox}
%     \end{minipage}%
%     \hfill
%     \begin{minipage}[t]{0.48\textwidth} % Adjust the width as needed
%         \centering
%         \begin{adjustbox}{max width=\textwidth}
%         \begin{tabular}{lcc}
%         \hline
%         \textbf{Statistically Insignificant} &  &  \\
%         \hline
%         PTA ID & Estimate & SE \\
%         \hline
%         641 & 2.028 & (1.255) \\
%         543 & 1.090 & (0.707) \\
%         509 & 0.210 & (0.216) \\
%         252 & 0.192 & (0.357) \\
%         508 & 0.140 & (0.122) \\
%         376 & 0.172 & (0.228) \\
%         416 & 0.424 & (0.295) \\
%         401 & 0.407 & (0.288) \\
%         152 & 0.110 & (0.266) \\
%         242 & 0.050 & (0.294) \\
%         390 & 0.0471 & (0.181) \\
%         396 & 0.019 & (0.379) \\
%         205 & 0.0012 & (0.178) \\
%         602 & -0.076 & (0.918) \\
%         383 & -0.202 & (0.152) \\
%         386 & -0.092 & (0.168) \\
%         84  & -0.059 & (0.120) \\
%         979  & -0.126 & (0.294) \\
%         644  & -0.189 & (0.122) \\
%         658  & -0.303 & (0.349) \\
%         \hline
%         \textbf{Negative and Statistically Significant} &  &  \\
%         \hline
%         PTA ID & Estimate & SE \\
%         \hline
%         399 & -0.473$^{\ast\ast\ast}$ & (0.127) \\
%         104 & -0.338$^{\ast\ast\ast}$ & (0.112) \\
%         677 & -1.366$^{\ast\ast\ast}$ & (0.385) \\
%         679 & -1.429$^{\ast\ast\ast}$ & (0.430) \\
%         323 & -0.338$^{\ast\ast}$ & (0.138) \\
%         512 & -0.458$^{\ast}$ & (0.266) \\
%         \hline
%         \end{tabular}
%         \end{adjustbox}
%     \end{minipage}
%     \vspace{0.5cm} % Optional: Add space between the tables and notes
%     \begin{adjustbox}{max width=\textwidth}
%     \begin{tabular}{l}
%     \hline
%     Exporter-Year FE & Yes \\
%     Importer-Year FE & Yes \\
%     Country-Pair FE & Yes \\
%     R-Squared & 0.998 \\
%     Observations & 73930 \\
%     \hline
%     \multicolumn{3}{l}{\footnotesize{Notes: Robust standard errors clustered at the country-pair level in parentheses.}} \\
%     \multicolumn{3}{l}{\footnotesize{Significance levels are indicated as follows: $^{\ast}$p$<$0.1; $^{\ast\ast}$p$<$0.05; $^{\ast\ast\ast}$p$<$0.01.}} \\
%     \end{tabular}
%     \end{adjustbox}
% \end{table}

%
\subsection{NS Estimation Results by Region}%
\label{subsec:NSEstimationResultsbyRegion}%
\subsubsection{NS Model Results by Region}%
\label{ssubsec:NSModelResultsbyRegion}%
\begin{table}[htbp]
    \centering
    \caption{Regional Results by TA Type}
    \label{tab:pta_types}
    \begin{adjustbox}{max width=\textwidth}
    \begin{tabular}{lccccc}
    \hline
     & \multicolumn{1}{c}{Africa} & \multicolumn{1}{c}{Americas} & \multicolumn{1}{c}{Asia} & \multicolumn{1}{c}{Europe} & \multicolumn{1}{c}{Intercontinental} \\
    \hline
    \textbf{Variables} &  &  &  &  &  \\
    \hline
    NN TA &  &  &  & 0.207$^{\ast\ast\ast}$ & 0.013 \\
     &  &  &  & (0.021) & (0.072) \\
    NN TA Lag &  &  &  & 0.192$^{\ast\ast\ast}$ & 0.016 \\
     &  &  &  & (0.023) & (0.073) \\
    NN TA + NN TA Lag &  &  &  & 0.399$^{\ast\ast\ast}$ & 0.029 \\
     &  &  &  & (0.026) & (0.102) \\
    \hline
    NS TA &  & 0.199$^{\ast\ast\ast}$ & -0.089 & 0.374$^{\ast\ast\ast}$ & 0.013 \\
     &  & (0.069) & (0.089) & (0.041) & (0.144) \\
    NS TA Lag &  & 0.234 & -0.067 & 0.349$^{\ast\ast\ast}$ & 0.231$^{\ast\ast\ast}$ \\
     &  & (0.190) & (0.060) & (0.041) & (0.061) \\
    NS TA + NS TA Lag &  & 0.434$^{\ast\ast}$ & -0.156$^{\ast}$ & 0.723$^{\ast\ast\ast}$ & 0.244 \\
     &  & (0.200) & (0.090) & (0.046) & (0.156) \\
    \hline
    SS TA & 0.578$^{\ast\ast\ast}$ & 0.476$^{\ast\ast\ast}$ & 0.153 & 0.530$^{\ast\ast\ast}$ & 0.004 \\
     & (0.154) & (0.139) & (0.117) & (0.107) & (0.121) \\
    SS TA Lag & -0.278 & -0.023 & -0.208$^{\ast\ast\ast}$ & 0.575$^{\ast\ast\ast}$ & 0.204$^{\ast\ast\ast}$ \\
     & (0.300) & (0.133) & (0.063) & (0.119) & (0.073) \\
    SS TA + SS TA Lag & 0.301 & 0.453$^{\ast\ast\ast}$ & -0.055 & 1.105$^{\ast\ast\ast}$ & 0.208 \\
     & (0.295) & (0.112) & (0.130) & (0.092) & (0.128) \\
    \hline
    Exporter-Year FE & Yes & Yes & Yes & Yes & Yes \\
    Importer-Year FE & Yes & Yes & Yes & Yes & Yes \\
    Country-Pair FE & Yes & Yes & Yes & Yes & Yes \\
    R-Squared & 0.997 & 0.999 & 0.999 & 0.997 & 0.998 \\
    Observations & 5838 & 10997 & 25308 & 28168 & 73930 \\
    \hline
    \multicolumn{6}{l}{\footnotesize{Notes: Robust standard errors clustered at the country-pair level in parentheses.}} \\
    \multicolumn{6}{l}{\footnotesize{Significance levels are indicated as follows: $^{\ast}$p$<$0.1; $^{\ast\ast}$p$<$0.05; $^{\ast\ast\ast}$p$<$0.01.}} \\
    \end{tabular}
    \end{adjustbox}
\end{table}


%
\subsection{NS PTA Estimation Results by Region}%
\label{subsec:NSPTAEstimationResultsbyRegion}%
\subsubsection{NS PTA Model Results by Region}%
\label{ssubsec:NSPTAModelResultsbyRegion}%
\begin{table}[htbp]
    \centering
    \caption{Africa PTA + PTA Lag Coefficients by Type}
    \label{tab:africa_pta}
    \begin{adjustbox}{max width=\textwidth}
    \begin{tabular}{lccc}
    \hline
    \textbf{PTA ID} & \textbf{NS PTA + Lag} & \textbf{SS PTA + Lag} & \textbf{NN PTA + Lag} \\
    \hline
    \textbf{NS and SS (or only NS)} &  &  &  \\
    \hline
    \multicolumn{4}{c}{No agreements in this category} \\
    \hline
    \textbf{Only SS} &  &  &  \\
    \hline
    670 &  & 0.326 &  \\
     &  & (0.410) &  \\
    787 &  & 0.304 &  \\
     &  & (0.233) &  \\
    \hline
    \textbf{Agreements with NN and NS} &  &  &  \\
    \hline
    \multicolumn{4}{c}{No agreements in this category} \\
    \hline
    Exporter-Year FE & Yes \\
    Importer-Year FE & Yes \\
    Country-Pair FE & Yes \\
    R-Squared & 0.997 \\
    Observations & 5838 \\
    \hline
    \multicolumn{4}{l}{\footnotesize{Notes: Robust standard errors clustered at the country-pair level in parentheses.}} \\
    \multicolumn{4}{l}{\footnotesize{Significance levels are indicated as follows: $^{\ast}$p$<$0.1; $^{\ast\ast}$p$<$0.05; $^{\ast\ast\ast}$p$<$0.01.}} \\
    \end{tabular}
    \end{adjustbox}
\end{table}
%
\begin{table}[htbp]
    \centering
    \caption{Americas PTA + PTA Lag Coefficients by Type}
    \label{tab:americas_pta}
    \begin{adjustbox}{max width=\textwidth}
    \begin{tabular}{lccc}
    \hline
    \textbf{PTA ID} & \textbf{NS PTA+Lag} & \textbf{SS PTA+Lag} & \textbf{NN PTA+Lag} \\
    \hline
    \textbf{Agreements with NS and SS (or only NS)} &  &  &  \\
    \hline
    188 & -0.811$^{\ast\ast\ast}$ & 0.685$^{\ast\ast}$ &  \\
    & (0.140) & (0.317) & \\
    163 & 0.346$^{\ast\ast\ast}$ &  &  \\
    & (0.098) &  &  \\
    168 & 0.410$^{\ast\ast\ast}$ &  &  \\
    & (0.113) &  &  \\
    218 & 0.879$^{\ast\ast\ast}$ &  &  \\
    & (0.172) &  &  \\
    645 & 0.117 &  &  \\
    & (0.141) &  &  \\
    \hline
    \textbf{Agreements with only SS} &  &  &  \\
    \hline
    141 &  & 0.265$^{\ast\ast\ast}$ &  \\
    &  & (0.024) &  \\
    213 &  & 1.342$^{\ast\ast\ast}$ &  \\
    &  & (0.435) &  \\
    239 &  & 0.572$^{\ast\ast\ast}$ &  \\
    &  & (0.173) &  \\
    616 &  & 0.488$^{\ast\ast\ast}$ &  \\
    &  & (0.044) &  \\
    201 &  & 0.545$^{\ast\ast}$ &  \\
    &  & (0.265) &  \\
    716 &  & 0.732$^{\ast\ast}$ &  \\
    &  & (0.358) &  \\
    612 &  & 0.517$^{\ast\ast}$ &  \\
    &  & (0.251) &  \\
    185 &  & 0.295 &  \\
    &  & (0.375) &  \\
    \hline
    \textbf{Agreements with NN and NS} &  &  &  \\
    \hline
    \multicolumn{4}{c}{No agreements in this category} \\
    \hline
    Exporter-Year FE & Yes \\
    Importer-Year FE & Yes \\
    Country-Pair FE & Yes \\
    R-Squared & 0.999 \\
    Observations & 10997 \\
    \hline
    \multicolumn{4}{l}{\footnotesize{Notes: Robust standard errors clustered at the country-pair level in parentheses.}} \\
    \multicolumn{4}{l}{\footnotesize{Significance levels are indicated as follows: $^{\ast}$p$<$0.1; $^{\ast\ast}$p$<$0.05; $^{\ast\ast\ast}$p$<$0.01.}} \\
    \end{tabular}
    \end{adjustbox}
\end{table}
%
\begin{table}[htbp]
    \centering
    \caption{Asia PTA + PTA Lag Coefficients by Type}
    \label{tab:asia_pta}
    \begin{adjustbox}{max width=\textwidth}
    \begin{tabular}{lccc}
    \hline
    \textbf{PTA ID} & \textbf{NS PTA + Lag} & \textbf{SS PTA + Lag} & \textbf{NN PTA + Lag} \\
    \hline
    \textbf{Agreements with NS and SS (or only NS)} &  &  &  \\
    \hline
    598 & 0.166$^{\ast\ast}$ &  &  \\
    & (0.083) &  &  \\
    71 & -0.138 &  &  \\
    &  & (0.091) &  \\
    \hline
    \textbf{Agreements with only SS} &  &  &  \\
    \hline
    70 &  & 0.472$^{\ast\ast\ast}$ &  \\
    &  & (0.150) &  \\
    100 &  & 0.376$^{\ast\ast\ast}$ &  \\
    &  & (0.105) &  \\
    67 &  & 0.342$^{\ast\ast\ast}$ &  \\
    &  & (0.125) &  \\
    683 &  & 1.080$^{\ast\ast\ast}$ &  \\
    &  & (0.237) &  \\
    599 &  & -0.967$^{\ast\ast\ast}$ &  \\
    &  & (0.191) &  \\
    220 &  & -1.215$^{\ast\ast\ast}$ &  \\
    &  & (0.093) &  \\
    221 &  & -2.955$^{\ast\ast\ast}$ &  \\
    &  & (0.0.727) &  \\
    675 &  & 1.360$^{\ast\ast}$ &  \\
    &  & (0.655) &  \\
    475 &  & 0.636$^{\ast\ast}$ &  \\
    &  & (0.298) &  \\
    1 &  & -0.732$^{\ast\ast}$ &  \\
    &  & (0.359) &  \\
    474 &  & 0.419$^{\ast}$ &  \\
    &  & (0.243) &  \\
    116 &  & 0.256 &  \\
    &  & (0.703) &  \\
    72 & 0.254 &  &  \\
    &  & (0.178) &  \\
    492 & 0.041 &  &  \\
    &  & (0.180) &  \\
    640 & 0.183 &  &  \\
    &  & (0.217) &  \\
    667 & -0.049 &  &  \\
    &  & (0.241) &  \\
    534 & -0.165 &  &  \\
    &  & (0.370) &  \\
    223 & -0.014 &  &  \\
    &  & (0.203) &  \\
    456 & -0.209 &  &  \\
    &  & (0.165) &  \\
    \hline
    \textbf{Agreements with NN and NS} &  &  &  \\
    \hline
    \multicolumn{4}{c}{No agreements in this category} \\
    \hline
    Exporter-Year FE & Yes \\
    Importer-Year FE & Yes \\
    Country-Pair FE & Yes \\
    R-Squared & 0.999 \\
    Observations & 25308 \\
    \hline
    \multicolumn{4}{l}{\footnotesize{Notes: Robust standard errors clustered at the country-pair level in parentheses.}} \\
    \multicolumn{4}{l}{\footnotesize{Significance levels are indicated as follows: $^{\ast}$p$<$0.1; $^{\ast\ast}$p$<$0.05; $^{\ast\ast\ast}$p$<$0.01.}} \\
    \end{tabular}
    \end{adjustbox}
\end{table}
%
\begin{center}
\small
\begin{longtable}{lccc}
    \caption{Europe PTA + PTA Lag Coefficients by Type} \label{tab:europe_pta} \\

    \hline
    \textbf{Agreements with NS} &  &  &  \\
    \hline
    \textbf{PTA ID} & \textbf{NS PTA+Lag} & \textbf{SS PTA+Lag} & \textbf{NN PTA+Lag} \\
    \hline
    \endfirsthead
    
    \multicolumn{4}{c}{{\bfseries \tablename\ \thetable{} -- continued from previous page}} \\
    \hline
    \textbf{PTA ID} & \textbf{NS PTA+Lag} & \textbf{SS PTA+Lag} & \textbf{NN PTA+Lag} \\
    \hline
    \endhead
    
    \hline \multicolumn{4}{r}{{Continued on next page}} \\ \hline
    \endfoot
    
    \hline
    \endlastfoot
    
    8 & 0.663$^{\ast\ast\ast}$ & 0.783$^{\ast\ast\ast}$ &  \\
    & (0.165) & (0.248) &  \\
    254 & 0.568$^{\ast\ast\ast}$ & 0.323 &  \\
    & (0.086) & (0.430) &  \\
    328 & 0.738$^{\ast\ast\ast}$ & 0.354 &  \\
    & (0.179) & (0.291) &  \\
    331 & 0.241 & -0.032 &  \\
    & (0.216) & (0.338) &  \\
    394 & 0.747$^{\ast\ast\ast}$ &  &  \\
    & (0.200) &  &  \\
    9 & 0.581$^{\ast\ast}$ &  &  \\
    & (0.285) &  &  \\
    255 & 0.171 &  &  \\
    & (0.235) &  &  \\
    389 & 0.411 &  &  \\
    & (0.333) &  &  \\
    \hline
    \textbf{Agreements with only SS} &  &  &  \\
    \hline
    \textbf{PTA ID} & \textbf{NS PTA+Lag} & \textbf{SS PTA+Lag} & \textbf{NN PTA+Lag} \\
    \hline
    5   &  & 3.811$^{\ast\ast\ast}$ &  \\
    &  & (0.278) &  \\
    132 &  & 2.241$^{\ast\ast\ast}$ &  \\
    &  & (0.252) &  \\
    7   &  & 1.153$^{\ast\ast\ast}$ &  \\
    &  & (0.271) &  \\
    13  &  & 2.303$^{\ast\ast\ast}$ &  \\
    &  & (0.246) &  \\
    128 &  & 2.773$^{\ast\ast\ast}$ &  \\
    &  & (0.213) &  \\
    129 &  & 0.556$^{\ast\ast\ast}$ &  \\
    &  & (0.208) &  \\
    135 &  & 0.696$^{\ast\ast\ast}$ &  \\
    &  & (0.206) &  \\
    192 &  & 1.199$^{\ast\ast\ast}$ &  \\
    &  & (0.152) &  \\
    621 &  & 0.614$^{\ast\ast\ast}$ &  \\
    &  & (0.186) &  \\
    133 &  & -0.707$^{\ast\ast\ast}$ &  \\
    &  & (0.225) &  \\
    11  &  & 0.663$^{\ast\ast}$ &  \\
    &  & (0.298) &  \\
    131  &  & 0.599$^{\ast\ast}$ &  \\
    &  & (0.276) &  \\
    154  &  & 0.773$^{\ast\ast}$ &  \\
    &  & (0.354) &  \\
    594  &  & 0.455$^{\ast}$ &  \\
    &  & (0.249) &  \\
    150 &  & 0.444 &  \\
    &  & (0.679) &  \\
    153 &  & 0.817 &  \\
    &  & (0.596) &  \\
    6   &  & 0.411 &  \\
    &  & (0.352) &  \\
    156   &  & -0.372 &  \\
    &  & (0.428) &  \\
    12   &  & -0.250 &  \\
    &  & (1.207) &  \\
    \hline
    \textbf{Agreements with NN} &  &  &  \\
    \hline
    \textbf{PTA ID} & \textbf{NS PTA+Lag} & \textbf{SS PTA+Lag} & \textbf{NN PTA+Lag} \\
    \hline
    335 & 0.727$^{\ast\ast\ast}$ & 1.099$^{\ast\ast\ast}$ & 0.399$^{\ast\ast\ast}$ \\
    & (0.047) & (0.098) & (0.025) \\
    \hline
    Exporter-Year FE & Yes \\
    Importer-Year FE & Yes \\
    Country-Pair FE & Yes \\
    R-Squared & 0.997 \\
    Observations & 28168 \\
    \hline
    \multicolumn{4}{l}{\footnotesize{Notes: Robust standard errors clustered at the country-pair level in parentheses.}} \\
    \multicolumn{4}{l}{\footnotesize{Significance levels are indicated as follows: $^{\ast}$p$<$0.1; $^{\ast\ast}$p$<$0.05; $^{\ast\ast\ast}$p$<$0.01.}} \\
\end{longtable}
\end{center}


% \begin{table}[htbp]
%     \centering
%     \caption{Europe PTA + PTA Lag Coefficients by Type}
%     \label{tab:europe_pta}
%     \begin{adjustbox}{max width=\textwidth}
%     \begin{tabular}{lccc}
%     \hline
%     \textbf{PTA ID} & \textbf{NS PTA+Lag} & \textbf{SS PTA+Lag} & \textbf{NN PTA+Lag} \\
%     \hline
%     \textbf{Agreements with NS} &  &  &  \\
%     \hline
%     8   & 0.663$^{\ast\ast\ast}$ & 0.783$^{\ast\ast\ast}$ &  \\
%     & (0.165) & (0.248) &  \\
%     254 & 0.568$^{\ast\ast\ast}$ & 0.323 &  \\
%     & (0.086) & (0.430) &  \\
%     328 & 0.738$^{\ast\ast\ast}$ & 0.354 &  \\
%     & (0.179) & (0.291) &  \\
%     394 & 0.747$^{\ast\ast\ast}$ &  &  \\
%     &  & (0.200) &  \\
%     9   & 0.581$^{\ast\ast}$ &  &  \\
%     & (0.285) &  &  \\
%     331   & 0.241 & -0.032 &  \\
%     & (0.216) & (0.338) &  \\
%     255   & 0.171 &  &  \\
%     & (0.235) &  &  \\
%     389   & 0.411 &  &  \\
%     & (0.333) &  &  \\
%     \hline
%     \textbf{Agreements with only SS} &  &  &  \\
%     \hline
%     5   &  & 3.811$^{\ast\ast\ast}$ &  \\
%     &  & (0.278) &  \\
%     132 &  & 2.241$^{\ast\ast\ast}$ &  \\
%     &  & (0.252) &  \\
%     7   &  & 1.153$^{\ast\ast\ast}$ &  \\
%     &  & (0.271) &  \\
%     13  &  & 2.303$^{\ast\ast\ast}$ &  \\
%     &  & (0.246) &  \\
%     128 &  & 2.773$^{\ast\ast\ast}$ &  \\
%     &  & (0.213) &  \\
%     129 &  & 0.556$^{\ast\ast\ast}$ &  \\
%     &  & (0.208) &  \\
%     135 &  & 0.696$^{\ast\ast\ast}$ &  \\
%     &  & (0.206) &  \\
%     192 &  & 1.199$^{\ast\ast\ast}$ &  \\
%     &  & (0.152) &  \\
%     621 &  & 0.614$^{\ast\ast\ast}$ &  \\
%     &  & (0.186) &  \\
%     133 &  & -0.707$^{\ast\ast\ast}$ &  \\
%     &  & (0.225) &  \\
%     11  &  & 0.663$^{\ast\ast}$ &  \\
%     &  & (0.298) &  \\
%     131  &  & 0.599$^{\ast\ast}$ &  \\
%     &  & (0.276) &  \\
%     154  &  & 0.773$^{\ast\ast}$ &  \\
%     &  & (0.354) &  \\
%     594  &  & 0.455$^{\ast}$ &  \\
%     &  & (0.249) &  \\
%     150 &  & 0.444 &  \\
%     &  & (0.679) &  \\
%     153 &  & 0.817 &  \\
%     &  & (0.596) &  \\
%     6   &  & 0.411 &  \\
%     &  & (0.352) &  \\
%     156   &  & -0.372 &  \\
%     &  & (0.428) &  \\
%     12   &  & -0.250 &  \\
%     &  & (1.207) &  \\
%     \hline
%     \textbf{Agreements with NN} &  &  &  \\
%     \hline
%     335 & 0.727$^{\ast\ast\ast}$ & 1.099$^{\ast\ast\ast}$ & 0.399$^{\ast\ast\ast}$ \\
%     & (0.047) & (0.098) & (0.025) \\
%     \hline
%     Exporter-Year FE & Yes \\
%     Importer-Year FE & Yes \\
%     Country-Pair FE & Yes \\
%     R-Squared & 0.997 \\
%     Observations & 28168 \\
%     \hline
%     \multicolumn{4}{l}{\footnotesize{Notes: Robust standard errors clustered at the country-pair level in parentheses.}} \\
%     \multicolumn{4}{l}{\footnotesize{Significance levels are indicated as follows: $^{\ast}$p$<$0.1; $^{\ast\ast}$p$<$0.05; $^{\ast\ast\ast}$p$<$0.01.}} \\
%     \end{tabular}
%     \end{adjustbox}
% \end{table}
%
\begin{center}
\small
\begin{longtable}{lccc}
    \caption{Intercont PTA+Lag Coefficients by Type} \label{tab:intercontinental_pta} \\

    \hline
    \textbf{Agreements with NS} &  &  &  \\
    \hline
    \textbf{PTA ID} & \textbf{NS PTA+Lag} & \textbf{SS PTA+Lag} & \textbf{NN PTA+Lag} \\
    \hline
    \endfirsthead
    
    \multicolumn{4}{c}{{\bfseries \tablename\ \thetable{} -- continued from previous page}} \\
    \hline
    \textbf{PTA ID} & \textbf{NS PTA+Lag} & \textbf{SS PTA+Lag} & \textbf{NN PTA+Lag} \\
    \hline
    \endhead
    
    \hline \multicolumn{4}{r}{{Continued on next page}} \\ \hline
    \endfoot
    
    \hline
    \endlastfoot
    
    17  & 0.800$^{\ast\ast\ast}$ & 1.055$^{\ast\ast\ast}$ &  \\
    & (0.249) & (0.244) &  \\
    330 & 0.286$^{\ast\ast\ast}$ & 0.662$^{\ast\ast\ast}$ &  \\
    & (0.088) & (0.188) &  \\
    304 & 0.787$^{\ast\ast\ast}$ & 0.591$^{\ast\ast}$ &  \\
    & (0.123) & (0.245) &  \\
    202 & 0.660$^{\ast\ast\ast}$ & 0.612$^{\ast\ast}$ &  \\
    & (0.124) & (0.241) &  \\
    679 & 0.546$^{\ast}$ & -1.636$^{\ast\ast\ast}$ &  \\
    & (0.310) & (0.400) &  \\
    323 & -0.335$^{\ast\ast}$ & -0.360 &  \\
    & (0.146) & (0.247) &  \\
    181 & 0.242 & 1.288$^{\ast\ast\ast}$ &  \\
    & (0.179) & (0.309) &  \\
    252 & 0.165 & 0.580 &  \\
    & (0.363) & (0.387) &  \\
    979 & -0.130 & 0.061 &  \\
    & (0.298) & (0.855) &  \\
    75  & 1.366$^{\ast\ast\ast}$ &  &  \\
    & (0.492) &  &  \\
    96  & 0.271$^{\ast\ast\ast}$ &  &  \\
    & (0.055) &  &  \\
    207 & 0.516$^{\ast\ast\ast}$ &  &  \\
    & (0.113) &  &  \\
    518 & 0.627$^{\ast\ast\ast}$ &  &  \\
    & (0.135) &  &  \\
    628 & 0.484$^{\ast\ast\ast}$ &  &  \\
    & (0.142) &  &  \\
    637 & 0.667$^{\ast\ast\ast}$ &  &  \\
    & (0.102) &  &  \\
    399 & -0.473$^{\ast\ast\ast}$ &  &  \\
    & (0.127) &  &  \\
    384 & 0.645$^{\ast}$ &  &  \\
    & (0.355) &  &  \\
    512 & -0.458$^{\ast}$ &  &  \\
    & (0.266) &  &  \\
    543 & 1.090 &  &  \\
    & (0.707) &  &  \\
    376 & 0.171 &  &  \\
    & (0.228) &  &  \\
    152 & 0.110 &  &  \\
    & (0.266) &  &  \\
    205 & 0.002 &  &  \\
    & (0.178) &  &  \\
    390 & 0.047 &  &  \\
    & (0.181) &  &  \\
    396 & 0.018 &  &  \\
    & (0.379) &  &  \\
    401 & 0.407 &  &  \\
    & (0.288) &  &  \\
    508 & 0.140 &  &  \\
    & (0.122) &  &  \\
    509 & 0.210 &  &  \\
    & (0.216) &  &  \\
    658 & -0.303 &  &  \\
    & (0.349) &  &  \\
    383 & -0.202 &  &  \\
    & (0.152) &  &  \\
    386 & -0.092 &  &  \\
    & (0.168) &  &  \\
    \hline
    \textbf{Agreements with only SS} &  &  &  \\
    \hline
    \textbf{PTA ID} & \textbf{NS PTA+Lag} & \textbf{SS PTA+Lag} & \textbf{NN PTA+Lag} \\
    \hline
    627 &  & 2.372$^{\ast\ast\ast}$ &  \\
    &  & (0.345) &  \\
    415   &  & 1.854$^{\ast\ast\ast}$ &  \\
    &  & (0.201) &  \\
    4   &  & 1.255$^{\ast\ast\ast}$ &  \\
    &  & (0.268) &  \\
    626 &  & 1.099$^{\ast\ast\ast}$ &  \\
    &  & (0.121) &  \\
    104 &  & -0.338$^{\ast\ast\ast}$ &  \\
    &  & (0.112) &  \\
    136 &  & 0.744$^{\ast\ast\ast}$ &  \\
    &  & (0.185) &  \\
    208 &  & 0.763$^{\ast\ast\ast}$ &  \\
    &  & (0.129) &  \\
    657 &  & 0.705$^{\ast\ast\ast}$ &  \\
    &  & (0.082) &  \\
    206 &  & 1.540$^{\ast\ast\ast}$ &  \\
    &  & (0.180) &  \\
    263 &  & 1.426$^{\ast\ast\ast}$ &  \\
    &  & (0.115) &  \\
    466 &  & 0.710$^{\ast\ast\ast}$ &  \\
    &  & (0.147) &  \\
    490 &  & 0.843$^{\ast\ast\ast}$ &  \\
    &  & (0.181) &  \\
    677 &  & -1.366$^{\ast\ast\ast}$ &  \\
    &  & (0.385) &  \\
    624 &  & 0.384$^{\ast\ast}$ &  \\
    &  & (0.163) &  \\
    15  &  & 0.313$^{\ast}$ &  \\
    &  & (0.179) &  \\
    227  &  & 0.348$^{\ast}$ &  \\
    &  & (0.186) &  \\
    242  &  & 0.050 &  \\
    &  & (0.294) &  \\
    416  &  & 0.424 &  \\
    &  & (0.295) &  \\
    641  &  & 2.027 &  \\
    &  & (1.255) &  \\
    644  &  & -0.190 &  \\
    &  & (0.122) &  \\
    602  &  & -0.076 &  \\
    &  & (0.918) &  \\
    \hline
    \textbf{Agreements with NN} &  &  &  \\
    \hline
    \textbf{PTA ID} & \textbf{NS PTA+Lag} & \textbf{SS PTA+Lag} & \textbf{NN PTA+Lag} \\
    \hline
    164 &  &  & 0.288$^{\ast\ast\ast}$ \\
    &  &  & (0.073) \\
    521 &  &  & 0.102$^{\ast\ast}$ \\
    &  &  & (0.045) \\
    84  &  &  & -0.059 \\
    &  &  & (0.120) \\
    \hline
    Exporter-Year FE & Yes \\
    Importer-Year FE & Yes \\
    Country-Pair FE & Yes \\
    R-Squared & 0.998 \\
    Observations & 73930 \\
    \hline
    \multicolumn{4}{l}{\footnotesize{Notes: Robust standard errors clustered at the country-pair level in parentheses.}} \\
    \multicolumn{4}{l}{\footnotesize{Significance levels are indicated as follows: $^{\ast}$p$<$0.1; $^{\ast\ast}$p$<$0.05; $^{\ast\ast\ast}$p$<$0.01.}} \\
\end{longtable}
\end{center}

% \begin{table}[htbp]
%     \centering
%     \caption{Intercontinental PTA + PTA Lag Coefficients by Type}
%     \label{tab:intercontinental_pta}
%     \begin{adjustbox}{max width=\textwidth}
%     \begin{tabular}{lccc}
%     \hline
%     \textbf{PTA ID} & \textbf{NS PTA + Lag} & \textbf{SS PTA + Lag} & \textbf{NN PTA + Lag} \\
%     \hline
%     \textbf{Agreements with NS and SS (or only NS)} &  &  &  \\
%     \hline
%     75  & 1.366$^{\ast\ast\ast}$ &  &  \\
%     & (0.492) &  &  \\
%     96  & 0.271$^{\ast\ast\ast}$ &  &  \\
%     & (0.055) &  &  \\
%     207 & 0.516$^{\ast\ast\ast}$ &  &  \\
%     & (0.113) &  &  \\
%     518 & 0.627$^{\ast\ast\ast}$ &  &  \\
%     & (0.135) &  &  \\
%     628 & 0.484$^{\ast\ast\ast}$ &  &  \\
%     & (0.142) &  &  \\
%     637 & 0.667$^{\ast\ast\ast}$ &  &  \\
%     & (0.102) &  &  \\
%     399 & -0.473$^{\ast\ast\ast}$ &  &  \\
%     & (0.127) &  &  \\
%     384 & 0.645$^{\ast}$ &  &  \\
%     & (0.355) &  &  \\
%     512 & -0.458$^{\ast}$ &  &  \\
%     & (0.266) &  &  \\
%     17  & 0.800$^{\ast\ast\ast}$ & 1.055$^{\ast\ast\ast}$ &  \\
%     & (0.249) & (0.244) &  \\
%     330 & 0.286$^{\ast\ast\ast}$ & 0.662$^{\ast\ast\ast}$ &  \\
%     & (0.088) & (0.188) &  \\
%     304 & 0.787$^{\ast\ast\ast}$ & 0.591$^{\ast\ast}$ &  \\
%     & (0.123) & (0.245) &  \\
%     202 & 0.660$^{\ast\ast\ast}$ & 0.612$^{\ast\ast}$ &  \\
%     & (0.124) & (0.241) &  \\
%     679 & 0.546$^{\ast}$ & -1.636$^{\ast\ast\ast}$ &  \\
%     & (0.310) & (0.400) &  \\
%     323 & -0.335$^{\ast\ast}$ & -0.360 &  \\
%     & (0.146) & (0.247) &  \\
%     181 & 0.242 & 1.288$^{\ast\ast\ast}$ &  \\
%     & (0.179) & (0.309) &  \\
%     252 & 0.165 & 0.580 &  \\
%     & (0.363) & (0.387) &  \\
%     979 & -0.130 & 0.061 &  \\
%     & (0.298) & (0.855) &  \\
%     543 & 1.090 &  &  \\
%     & (0.707) &  &  \\
%     376 & 0.171 &  &  \\
%     & (0.228) &  &  \\
%     152 & 0.110 &  &  \\
%     & (0.266) &  &  \\
%     205 & 0.002 &  &  \\
%     & (0.178) &  &  \\
%     390 & 0.047 &  &  \\
%     & (0.181) &  &  \\
%     396 & 0.018 &  &  \\
%     & (0.379) &  &  \\
%     401 & 0.407 &  &  \\
%     & (0.288) &  &  \\
%     508 & 0.140 &  &  \\
%     & (0.122) &  &  \\
%     509 & 0.210 &  &  \\
%     & (0.216) &  &  \\
%     658 & -0.303 &  &  \\
%     & (0.349) &  &  \\
%     383 & -0.202 &  &  \\
%     & (0.152) &  &  \\
%     386 & -0.092 &  &  \\
%     & (0.168) &  &  \\
%     \hline
%     \textbf{Agreements with only SS} &  &  &  \\
%     \hline
%     627 &  & 2.372$^{\ast\ast\ast}$ &  \\
%     &  & (0.345) &  \\
%     415   &  & 1.854$^{\ast\ast\ast}$ &  \\
%     &  & (0.201) &  \\
%     4   &  & 1.255$^{\ast\ast\ast}$ &  \\
%     &  & (0.268) &  \\
%     626 &  & 1.099$^{\ast\ast\ast}$ &  \\
%     &  & (0.121) &  \\
%     104 &  & -0.338$^{\ast\ast\ast}$ &  \\
%     &  & (0.112) &  \\
%     136 &  & 0.744$^{\ast\ast\ast}$ &  \\
%     &  & (0.185) &  \\
%     208 &  & 0.763$^{\ast\ast\ast}$ &  \\
%     &  & (0.129) &  \\
%     657 &  & 0.705$^{\ast\ast\ast}$ &  \\
%     &  & (0.082) &  \\
%     206 &  & 1.540$^{\ast\ast\ast}$ &  \\
%     &  & (0.180) &  \\
%     263 &  & 1.426$^{\ast\ast\ast}$ &  \\
%     &  & (0.115) &  \\
%     466 &  & 0.710$^{\ast\ast\ast}$ &  \\
%     &  & (0.147) &  \\
%     490 &  & 0.843$^{\ast\ast\ast}$ &  \\
%     &  & (0.181) &  \\
%     677 &  & -1.366$^{\ast\ast\ast}$ &  \\
%     &  & (0.385) &  \\
%     624 &  & 0.384$^{\ast\ast}$ &  \\
%     &  & (0.163) &  \\
%     15  &  & 0.313$^{\ast}$ &  \\
%     &  & (0.179) & () \\
%     227  &  & 0.348$^{\ast}$ &  \\
%     &  & (0.186) &  \\
%     242  &  & 0.050 &  \\
%     &  & (0.294) &  \\
%     416  &  & 0.424 &  \\
%     &  & (0.295) &  \\
%     641  &  & 2.027 &  \\
%     &  & (1.255) &  \\
%     644  &  & -0.190 &  \\
%     &  & (0.122) &  \\
%     602  &  & -0.076 &  \\
%     &  & (0.918) &  \\
%     \hline
%     \textbf{Agreements with NN and NS} &  &  &  \\
%     \hline
%     164 &  &  & 0.288$^{\ast\ast\ast}$ \\
%     &  &  & (0.073) \\
%     521 &  &  & 0.102$^{\ast\ast}$ \\
%     &  &  & (0.045) \\
%     84  &  &  & -0.059 \\
%     &  &  & (0.120) \\
%     \hline
%     Exporter-Year FE & Yes \\
%     Importer-Year FE & Yes \\
%     Country-Pair FE & Yes \\
%     R-Squared & 0.998 \\
%     Observations & 73930 \\
%     \hline
%     \multicolumn{4}{l}{\footnotesize{Notes: Robust standard errors clustered at the country-pair level in parentheses.}} \\
%     \multicolumn{4}{l}{\footnotesize{Significance levels are indicated as follows: $^{\ast}$p$<$0.1; $^{\ast\ast}$p$<$0.05; $^{\ast\ast\ast}$p$<$0.01.}} \\
%     \end{tabular}
%     \end{adjustbox}
% \end{table}

%
\newpage%
\subsubsection{Additional Findings}%
\label{ssubsec:AdditionalFindings}%
Here you can include additional findings and discussions.

%
\section{Analysis and Discussion}%
\label{sec:AnalysisandDiscussion}%
This is the analysis and discussion

%
\section{Conclusion}%
\label{sec:Conclusion}%
This is the conclusion

%
\newpage%
%TC:ignore%
\section{References}%
\label{sec:References}%
\printbibliography

%
%TC:endignore%
\newpage%
%TC:ignore%
\section{Appendix}%
\label{sec:Appendix}%
\subsection{Subsection in Appendix}%
\label{subsec:SubsectioninAppendix}%
Content in the appendix should not be counted in the word count.

%
%TC:endignore%
\end{document}